\chapter{Conservación genealógica de la información, entropías de reducción y estructura térmica}
\label{chap:informacion-entropia-landauer}

La información que transporta un estado HMT no se agota en la coordenada que
publica uno de sus lectores. La coordenada visible, la hoja, la orientación,
el cociente, el acarreo, la ruta y el sello de cierre pertenecen al mismo
estado enriquecido, aunque una evaluación concreta sólo muestre una parte de
él. Esta distinción permite situar en un único marco tres operaciones que la
descripción macroscópica suele presentar por separado: conservar una historia,
reducirla a una lectura y atribuir energía a la información efectivamente
borrada.

El capítulo desarrolla esa cadena sin identificar sus escalones. Primero se
formula la memoria de la fibra y el criterio exacto de recuperabilidad. Después
se demuestra la nulidad del término lógico de Landauer para una actualización
biyectiva del estado enriquecido de TPK. A continuación se distinguen cinco
entropías que poseen dominios y funciones diferentes. Por último, el cociente
areal--volumétrico y el reloj de \(108\) pasos proporcionan la estructura que
enlaza información triádica, acción y temperatura.

\section{Desintegración de la lectura visible sobre fibras y recuperación de la memoria genealógica}

Sea \(\Omega\) un espacio finito de estados enriquecidos y sea
\[
 q:\Omega\longrightarrow Y
\]
un lector visible. El valor \(q(x)\) es una coordenada del estado \(x\), no su
descripción exhaustiva. Denotemos por
\[
 M:\Omega\longrightarrow\mathcal M
\]
el registro conjunto de las coordenadas genealógicas que el lector no publica:
hoja, orientación, cociente o acarreo, fase, ruta y sello, según el dominio
considerado.

\begin{definition}[Lector enriquecido y memoria de fibra]
El lector enriquecido asociado a \(q\) y \(M\) es
\begin{equation}
 \widehat q=(q,M):\Omega\longrightarrow Y\times\mathcal M.
 \label{eq:lector-enriquecido-informacion}
\end{equation}
La memoria de fibra de \(q\) es la información necesaria para distinguir los
estados contenidos en una misma fibra \(q^{-1}(y)\).
\end{definition}

La extensión residuo--cociente de APP es el ejemplo aritmético elemental. En
cada celda se conservan las dos divisiones euclídeas
\[
 i+j=9q^+_{ij}+R^+_{ij},
 \qquad
 ij=9q^\times_{ij}+R^\times_{ij}.
\]
La reducción módulo nueve publica \(R^+\) o \(R^\times\); el levantamiento
conserva simultáneamente \((R,q)\). Al sumar sobre el dominio nonádico se
obtiene la identidad exacta
\begin{equation}
 2025-810
 =54+9\cdot129.
 \label{eq:app-residuo-cociente-informacion}
\end{equation}
El primer término de la derecha es la diferencia entre los residuos visibles;
el segundo recompone la diferencia entre los cocientes. La igualdad demuestra
que la proyección residual no determina por sí sola la operación entera.

Sea ahora \(X\) una variable aleatoria con valores en \(\Omega\). Trabajaremos
con entropías adimensionales en nats. La regla de la cadena da
\begin{equation}
 \mathsf H(X)
 =\mathsf H(q(X))+\mathsf H(X\mid q(X)).
 \label{eq:cadena-entropia-fibra}
\end{equation}
El segundo término cuantifica la información genealógica que permanece dentro
de la fibra de la lectura.

\begin{theorem}[Recuperabilidad por memoria de fibra]
\label{thm:recuperabilidad-memoria-fibra}
Si el lector enriquecido \(\widehat q=(q,M)\) es inyectivo sobre el soporte de
\(X\), entonces
\begin{equation}
 \boxed{\mathsf H(X\mid q(X),M(X))=0.}
 \label{eq:entropia-condicional-memoria-cero}
\end{equation}
La entropía \(\mathsf H(X\mid q(X))\) mide exactamente la información que se
perdería al conservar la coordenada visible y descartar la memoria.
\end{theorem}

\begin{proof}
La inyectividad de \(\widehat q\) proporciona una inversa sobre su imagen. Por
tanto, el par \((q(X),M(X))\) determina \(X\) y la entropía condicional es
cero. Si se omite \(M(X)\), las entradas de una misma fibra quedan
identificadas; la regla de la cadena \eqref{eq:cadena-entropia-fibra} asigna a
esa identificación el término \(\mathsf H(X\mid q(X))\).
\end{proof}

La definición genealógica de número se inserta precisamente aquí. Un número
HMT es una sección compatible de cilindros supervivientes junto con la memoria
de su ruta, orientación, firma y cierre. Su valor real es la evaluación de esa
sección. La evaluación puede olvidar coordenadas; la sección las conserva.

\section{Reversibilidad lógica y término de Landauer}

Sea
\[
 U:\Omega\longrightarrow\Omega
\]
una actualización del estado enriquecido de TPK. La pregunta de Landauer se
formula sobre el mapa completo, antes de aplicar un lector visible o reiniciar
un registro; su formulación termodinámica del borrado lógico se remonta a
Landauer \cite{Landauer1961}.

\begin{theorem}[Landauer nulo para el transporte completo]
\label{thm:landauer-nulo-tpk-anexo}
Si \(U\) es biyectiva, entonces, para toda distribución de entrada \(X\),
\begin{equation}
 \mathsf H(X\mid U(X))=0.
 \label{eq:landauer-entropia-cero-anexo}
\end{equation}
En el régimen ideal de Landauer, el término mínimo asociado exclusivamente al
borrado lógico interno es
\begin{equation}
 \boxed{
 Q_{\rm L}^{\min}[U]
 =k_B\Theta\,\mathsf H(X\mid U(X))=0.}
 \label{eq:landauer-nulo-anexo}
\end{equation}
\end{theorem}

\begin{proof}
La biyección proporciona \(U^{-1}\), de modo que
\(X=U^{-1}(U(X))\). La salida completa determina la entrada y la entropía
condicional se anula. La forma informacional del principio de Landauer asigna
el coste de borrado al término \(k_B\Theta\mathsf H(X\mid U(X))\), que por la
igualdad anterior vale cero.
\end{proof}

El resultado localiza la conservación en el espacio enriquecido. Para una
salida proyectada \(q(U(X))\), la información no publicada es
\begin{equation}
 \mathsf H\bigl(X\mid q(U(X))\bigr).
 \label{eq:informacion-proyeccion-anexo}
\end{equation}
Un reinicio físico que identifique estados distintos debe transferir esa
información al entorno. En particular, el reset de un TRIT uniforme a un único
estado posee
\[
 \mathsf H(X)=\log3,
 \qquad
 \boxed{Q_{\rm reset}\geq k_B\Theta\log3.}
\]
La actualización reversible, la proyección y el reinicio son así operaciones
distintas. La primera conserva la genealogía; la segunda reduce la lectura; la
tercera elimina físicamente un registro. La disipación adicional debida al
control, al ruido o a una velocidad finita pertenece a la realización del
dispositivo y no al término lógico de borrado de
\eqref{eq:landauer-nulo-anexo}.

\section{\(5\) nociones de entropía y sus dominios}

La palabra \emph{entropía} designa en esta construcción cinco funcionales
relacionados, pero no intercambiables. Su separación evita atribuir temperatura
a una cuenta puramente combinatoria o pérdida de información a una evolución
que sólo cambia de lectura.

\subsection{Entropía condicional de las fibras del lector visible y pérdida de genealogía}

La cantidad
\[
 \mathsf H_{\rm fib}(q;X):=\mathsf H(X\mid q(X))
\]
mide la información de la genealogía que una lectura visible no distingue.
Su dominio es un lector y una distribución sobre el estado enriquecido. Se
anula cuando el lector es fiel sobre el soporte considerado.

\subsection{Entropía informacional de una decisión triádica}

Para una distribución \(p=(p_{-1},p_0,p_{+1})\) sobre los tres estados
visibles del TRIT,
\[
 I_{\rm TRIT}(p)=-\sum_{a\in\{-1,0,+1\}}p_a\log p_a.
\]
La distribución uniforme produce
\begin{equation}
 \boxed{I_{\rm TRIT}=\log3.}
 \label{eq:informacion-trit-log3-anexo}
\end{equation}
Esta cantidad mide una decisión local entre tres alternativas visibles.

\subsection{Entropía topológica de las historias admisibles}

Si \(A\) es la matriz de adyacencia de un subshift finito, su entropía
topológica es
\[
 h_{\rm top}(A)=\log\rho(A),
\]
donde \(\rho(A)\) es el radio espectral. Este funcional mide la tasa
asintótica de crecimiento de las historias compatibles, no la incertidumbre
de una decisión aislada.

\subsection{Entropía combinatoria y máxima entropía estadística}

Para niveles de energía \(\varepsilon_i\) con degeneraciones \(g_i\), una
ocupación fermiónica \(n_i\in\{0,\ldots,g_i\}\) posee
\[
 \Omega(g_i,n_i)=\binom{g_i}{n_i}
\]
microestados. Con número y energía totales fijados,
\begin{equation}
 \Omega(N,E)=
 \sum_{\substack{(n_i):\ \sum_i n_i=N\\
                         \sum_i\varepsilon_i n_i=E}}
 \prod_i\binom{g_i}{n_i},
 \qquad
 S_{\rm mc}=\log\Omega(N,E).
 \label{eq:entropia-microcanonica-anexo}
\end{equation}
Al fijar número y energía como valores medios, la maximización de Shannon
\cite{Shannon1948} da
la distribución producto
\[
 \mathbb P(n_1,\ldots,n_M)
 \propto
 \exp\!\left[-\beta\sum_j\varepsilon_jn_j
              +\beta\mu\sum_j n_j\right]
\]
y, por tanto,
\begin{equation}
 \boxed{
 \langle n_j\rangle_{\rm FD}
 =\frac{1}{e^{\beta(\varepsilon_j-\mu)}+1}.}
 \label{eq:fermi-dirac-entropia-anexo}
\end{equation}
La ocupación bosónica sustituye el conteo binomial por
\(\binom{g+n-1}{n}\) y produce
\[
 \langle n_j\rangle_{\rm BE}
 =\frac{1}{e^{\beta(\varepsilon_j-\mu)}-1}.
\]
En ambos casos, \(\beta\) y \(\mu\) quedan determinados por las restricciones
macroscópicas declaradas.

\subsection{Entropía de von Neumann de \(1\) estado reducido}

Para una matriz de densidad \(\rho_A\) sobre una región o subálgebra,
\[
 S_{\rm vN}(\rho_A)=-\operatorname{Tr}(\rho_A\log\rho_A).
\]
Su dominio es un estado reducido. Cuando \(\rho_A\) procede de una
purificación global, esta cantidad mide el entrelazamiento a través del corte.
El capítulo siguiente construye la reducción triádica, el Hamiltoniano modular
y la cuantización de área que le dan una residencia geométrica precisa.

\begin{proposition}[Separación tipada de las cinco entropías]
\label{prop:cinco-entropias-tipadas}
Las cantidades \(\mathsf H_{\rm fib}\), \(I_{\rm TRIT}\), \(h_{\rm top}\),
\(S_{\rm mc}\) y \(S_{\rm vN}\) no se sustituyen entre sí: dependen,
respectivamente, de una fibra de lectura, una distribución local triádica, un
sistema de historias, un conjunto de microestados compatible con restricciones
y un estado reducido. Un mapa físico puede relacionarlas sólo después de
declarar el coarse--graining, el ensamble, el Hamiltoniano o la purificación
que cambia de un dominio al siguiente.
\end{proposition}

\begin{proof}
Cada funcional está definido sobre un tipo de entrada diferente. La entropía
condicional necesita un lector; \(I_{\rm TRIT}\), una distribución sobre tres
símbolos; \(h_{\rm top}\), una dinámica de palabras; \(S_{\rm mc}\), una
restricción macrocanónica; y \(S_{\rm vN}\), un operador de densidad. Una
igualdad numérica accidental no suministra los mapas que cambian esos tipos.
\end{proof}

\section{Incidencia areal--volumétrica y entropía topológica}

El calendario de TPK induce la bisagra entre los modos aditivo y
multiplicativo. En el bloque primitivo
\[
 \tau_3=(+1,-1,0),
\]
la regla
\[
 +1:m\mapsto\Sigma,
 \qquad
 -1:m\mapsto\Pi,
 \qquad
 0:m\mapsto m
\]
produce la órbita cíclica \((\Sigma,\Pi,\Pi)\). El lector geométrico
\[
 \Sigma\longmapsto\mathscr H_{\rm ar},
 \qquad
 \Pi\longmapsto\mathscr H_{\rm vol}
\]
transporta esos modos a las hojas areal y volumétrica.

\begin{theorem}[Incidencia necesaria de las hojas]
\label{thm:incidencia-hojas-anexo}
En el orden \((\mathscr H_{\rm ar},\mathscr H_{\rm vol})\), la matriz
primitiva de incidencias es
\begin{equation}
 \boxed{
 F_{\rm av}=
 \begin{pmatrix}0&1\\1&1\end{pmatrix}.}
 \label{eq:fav-anexo-informacion}
\end{equation}
Los calendarios de longitudes \(9,27,108\) producen
\[
 N_9=3F_{\rm av},
 \qquad N_{27}=9F_{\rm av},
 \qquad N_{108}=36F_{\rm av}.
\]
\end{theorem}

\begin{proof}
Los pares consecutivos del ciclo son
\(\Sigma\to\Pi\), \(\Pi\to\Pi\) y \(\Pi\to\Sigma\). Cada uno aparece una
vez y \(\Sigma\to\Sigma\) no aparece. Esto fija las cuatro entradas de
\(F_{\rm av}\). Los calendarios largos contienen, respectivamente, tres,
nueve y treinta y seis copias del bloque primitivo.
\end{proof}

\begin{theorem}[Ley de Fibonacci y medida de máxima entropía]
\label{thm:fibonacci-parry-anexo}
Para \(n\geq1\),
\begin{equation}
 F_{\rm av}^{n}
 =\begin{pmatrix}F_{n-1}&F_n\\F_n&F_{n+1}\end{pmatrix}.
 \label{eq:potencias-fav-anexo}
\end{equation}
El subshift de memoria uno definido por \(F_{\rm av}\) tiene
\begin{equation}
 \boxed{h_{\rm top}=\log\varphi},
 \qquad
 \boxed{\frac{\mu_{\rm ar}}{\mu_{\rm vol}}=\varphi^{-2}}.
 \label{eq:entropia-parry-hojas-anexo}
\end{equation}
\end{theorem}

\begin{proof}
La identidad matricial se obtiene por inducción a partir de
\(F_{n+1}=F_n+F_{n-1}\). La raíz de Perron de \(F_{\rm av}\) es \(\varphi\),
de donde \(h_{\rm top}=\log\varphi\). Como la matriz es simétrica, los pesos
de Parry \cite{Parry1964} son proporcionales al cuadrado de un vector de Perron
\((1,\varphi)\), es decir, a \((1,\varphi^2)\).
\end{proof}

Las expresiones \(\log3\) y \(\log\varphi\) quedan así separadas por su
genealogía. La primera cuantifica una elección local uniforme entre tres
estados; la segunda mide el crecimiento de las historias permitidas por la
incidencia de hojas. La relación areal--volumétrica es una propiedad del
retorno completo y por ello aparece con el exponente dos.

\section{Transducción de Boltzmann y energía por decisión triádica}

Sea \(\mathcal L_T\) la recta de duración, \(\mathcal L_\Theta\) la recta de
temperatura, \(\mathcal L_S\) la recta de acción y \(\mathcal L_E\) la recta
de energía. La constante de Boltzmann es el isomorfismo positivo
\begin{equation}
 \boxed{
 k_B:\mathcal L_\Theta\longrightarrow\mathcal L_E,
 \qquad [k_B]=\mathbf E-\boldsymbol\Theta.}
 \label{eq:boltzmann-transductor-anexo}
\end{equation}
Así, \(\beta=(k_B\Theta)^{-1}\) pertenece a la recta inversa de energía.

El lector determinantal de la acción construye
\[
 \hbar_{\rm int}
 =\eta_{\rm ret}\,10^{-34}\mathcal U_S\in\mathcal L_S.
\]
Si \(t_0\in\mathcal L_T^+\) es la duración interna elemental, el reloj de
retorno completo fija
\[
 \omega_0=\frac{2\pi}{108t_0},
 \qquad
 \varepsilon_{\rm clk}=\hbar_{\rm int}\omega_0\in\mathcal L_E.
\]

\begin{definition}[Escala térmica triádica]
La sección \(\Theta_{\rm clk}\in\mathcal L_\Theta^+\) se define por
\begin{equation}
 \boxed{
 k_B\Theta_{\rm clk}\log3
 =\hbar_{\rm int}\frac{2\pi}{108t_0}.}
 \label{eq:boltzmann-accion-informacion-anexo}
\end{equation}
\end{definition}

La igualdad fija la energía por decisión triádica y enlaza tres objetos ya
construidos: la información local \(\log3\), la acción interna y la frecuencia
del reloj. El producto \(k_B\Theta_{\rm clk}\) queda determinado dentro de la
carta interna. La elección de bases térmica y energética separa después la
coordenada de \(k_B\) y la de \(\Theta_{\rm clk}\). En la carta kelvin--joule,
la definición metrológica publica
\[
 k_B^{\rm SI}=1.380649\times10^{-23}\ \mathrm{J\,K^{-1}}.
\]

\section{Operador térmico sobre las historias}

Sea \(\mathcal E\) el conjunto de aristas de un grafo finito de extensiones,
sea \(a_*(e)\geq0\) el incremento adimensional de acción y sea
\(\epsilon_*\in\mathcal L_E^+\) una sección positiva de energía. Para
\(\beta\in\mathcal L_E^{-1}\), definimos
\begin{equation}
 (\mathcal L_{\beta,*}f)(y)
 =\sum_{e:x\to y}
 \exp[-\beta\epsilon_*a_*(e)]f(x).
 \label{eq:operador-transferencia-termica-anexo}
\end{equation}
El exponente es adimensional y la positividad del operador conserva la
composición de las historias. La normalización de presión selecciona el estado
térmico mediante
\begin{equation}
 \boxed{\rho(\mathcal L_{\beta,*})=1.}
 \label{eq:presion-espectral-termica-anexo}
\end{equation}

\begin{proposition}[Confluencia de la rama informacional y la rama espectral]
\label{prop:confluencia-termica-anexo}
Una realización térmica HMT--MD queda determinada cuando un mismo valor de
\(\beta=(k_B\Theta)^{-1}\) satisface la normalización espectral
\eqref{eq:presion-espectral-termica-anexo} y la escala acción--información
\eqref{eq:boltzmann-accion-informacion-anexo}. La primera ecuación pondera las
historias; la segunda fija la energía de la decisión local. Su confluencia
convierte la contabilidad genealógica en un ensamble térmico sin identificar
información, entropía y energía.
\end{proposition}

La termodinámica aparece, por tanto, después de la conservación genealógica.
El estado completo mantiene la historia; el lector define el coarse--graining;
el conteo organiza los microestados; y el transductor de Boltzmann asigna una
escala energética a la información bajo una temperatura declarada.
